\begin{exsol}{17}{\MR{} Mathematical Exercise~2.2}{53}{Exam}{B1}
\asked In the beach-improvement setting of Q~2.1, show that equal cost sharing makes one family
reject the efficient project, find the largest project all families accept under equal shares,
and exhibit a different spending level and cost split that every family prefers.

\method Participation (voluntary agreement) constraints under a fixed sharing rule; the value
maximization principle and its second sentence: an inefficient allocation is strictly
Pareto dominated by a value-maximizing one once transfers are free (\MR{} p.~36; the construction
in Q~2.3, item~26).

\solution
Values as in item~16: $v_1 = v_2 = 5y - \tfrac12 y^2$, $v_3 = 7y - \tfrac12 y^2$,
$v_4 = 4y - y^2$; total value $W(y) = 20y - \tfrac52 y^2$, maximized at $y^* = 4$ with $W(4) = 40$.
Each family can refuse, in which case nothing is built and it gets $0$.

\textbf{(a) Equal shares at the efficient level.} With equal shares each family pays $y/4$, so its
net value is $n_i(y) = v_i(y) - y/4$:
\[
  n_1(y) = n_2(y) = 4.75\,y - \tfrac12 y^2, \qquad n_3(y) = 6.75\,y - \tfrac12 y^2, \qquad
  n_4(y) = 3.75\,y - y^2 .
\]
At $y^* = 4$: $n_1 = n_2 = 19 - 8 = 11$, $n_3 = 27 - 8 = 19$, and
\[
  n_4(4) = 15 - 16 = -1 < 0 .
\]
Family~4 values the improvement at $v_4(4) = 0$ but would pay $1$: it is better off with no
project and will not agree. (The net values still sum to $W(4) = 11 + 11 + 19 - 1 = 40$; the problem
is the division, not the size of the pie.)

\textbf{(b) Largest unanimous project under equal shares.} Each family accepts $y$ iff
$n_i(y) \ge 0$:
\[
  n_1, n_2 \ge 0 \iff y \le 9.5, \qquad n_3 \ge 0 \iff y \le 13.5, \qquad
  n_4 = y(3.75 - y) \ge 0 \iff y \le 3.75 .
\]
The binding constraint is family~4's, so the largest project everyone accepts is
\[
  \bar y = 3.75 ,
\]
at which family~4 is just indifferent (any $y < 3.75$ it strictly prefers to nothing). Net values
at $\bar y$: $n_1 = n_2 = 10.78125$, $n_3 = 18.28125$, $n_4 = 0$, so
$W(3.75) = 39.84375$. (Family~4's \emph{favourite} level under equal shares is $1.875$, where
$n_4' = 3.75 - 2y = 0$; the question asks for the largest level it would still accept.)

\textbf{(c) $\bar y$ is inefficiently low.} $\bar y = 3.75 < 4 = y^*$ and
$W(4) - W(3.75) = 40 - 39.84375 = 0.15625 = 5/32 > 0$. Since total value is higher at $y^* = 4$,
there must be a cost split at $y = 4$ that every family prefers. Following the construction of
Q~2.3, give each family its net value at $\bar y$ plus an equal share of the gain, $5/128$
each, and round to cents. Family $i$ pays $t_i = v_i(4) - n_i(\bar y) - 5/128$:
\begin{center}
\begin{tabular}{lccccc}
\toprule
Family & 1 & 2 & 3 & 4 & Total \\
\midrule
Value $v_i(4)$ & $12$ & $12$ & $20$ & $0$ & $44$ \\
Contribution $t_i$ & $1.18$ & $1.18$ & $1.68$ & $-0.04$ & $4.00$ \\
Net value at $y=4$ & $10.82$ & $10.82$ & $18.32$ & $0.04$ & $40.00$ \\
Net value at $\bar y = 3.75$, equal shares & $10.78125$ & $10.78125$ & $18.28125$ & $0$ & $39.84375$ \\
\bottomrule
\end{tabular}
\end{center}
The contributions cover the cost exactly ($1.18 + 1.18 + 1.68 - 0.04 = 4$), and every family is
strictly better off than under the equal-share plan: the three families with high marginal values
spend a little more and pay family~4 a side payment of $\$0.04$ to accept the larger project.

\emph{Without a payment to family~4.} If contributions must be non-negative, a level just below
$y^*$ works because family~4 then values the project positively. At $y = 3.9$
($v = 11.895,\ 11.895,\ 19.695,\ 0.39$), contributions $(1.10, 1.10, 1.40, 0.30)$ sum to $3.9$
and give net values $(10.795,\ 10.795,\ 18.295,\ 0.09)$, again strictly above
$(10.78125, 10.78125, 18.28125, 0)$.

\begin{center}
\begin{tikzpicture}
\begin{axis}[width=0.66\linewidth, height=5.4cm, xmin=0, xmax=5.2, ymin=-7, ymax=23,
  xlabel={total spending $y$}, ylabel={net value with equal shares}, axis lines=left,
  legend style={font=\footnotesize, draw=none, at={(0.02,0.98)}, anchor=north west},
  tick label style={font=\footnotesize}, clip=false]
  \addplot[thick, hdr, domain=0:5] {6.75*x-0.5*x^2}; \addlegendentry{family 3}
  \addplot[thick, ans, domain=0:5] {4.75*x-0.5*x^2}; \addlegendentry{families 1, 2}
  \addplot[thick, red!70!black, domain=0:5] {3.75*x-x^2}; \addlegendentry{family 4}
  \draw[gray] (axis cs:0,0) -- (axis cs:5.2,0);
  \draw[dotted] (axis cs:3.75,-7) -- (axis cs:3.75,23);
  \draw[dotted] (axis cs:4,-7) -- (axis cs:4,23);
  \node[font=\footnotesize, anchor=south east] at (axis cs:3.75,23) {$\bar y$};
  \node[font=\footnotesize, anchor=south west] at (axis cs:4,23) {$y^*$};
\end{axis}
\end{tikzpicture}
\end{center}

\interp A fixed cost-sharing rule ties the division ($x$) to the decision ($y$), which destroys
the separation behind the value maximization principle: the low-value family's veto caps the
project below the efficient level. Side payments (or any flexible cost split) restore the
separation, and then unanimity and efficiency coincide. This is the organizational lesson of
\MR{} pp.~36--38: rules that restrict transfers (equal shares, uniform prices, one-size-fits-all
policies) can block value-creating decisions, so the question is always whether the transfer
needed is feasible and affordable.

\verify $n_4(3.75) = 14.0625 - 14.0625 = 0$; $n_4(4) = -1$. $W(3.75) = 75 - 35.15625 = 39.84375$.
The contributions in the table sum to $4$ and each net value exceeds its benchmark; likewise at
$y = 3.9$ (\texttt{check\_C.py} checks all entries exactly in rationals).

\answer{With equal shares family~4's net value at $y^* = 4$ is $4\cdot4 - 16 - 1 = -1 < 0$, so it
refuses. The largest project all accept is $\bar y = 3.75$ (family~4's
$3.75y - y^2 \ge 0$), with total value $39.84375 < 40$. Spending $4$ with contributions
$(1.18, 1.18, 1.68, -0.04)$ gives net values $(10.82, 10.82, 18.32, 0.04)$, strictly above
$(10.78125, 10.78125, 18.28125, 0)$ for every family.}
\end{exsol}
